\section{Metrics \& Experimental Setup}

\subsection{Utility-Oriented Multi-Dimensional Metrics}

\eva{} proposes a multi-dimensional evaluation metric system that moves beyond simple accuracy scores toward utility-oriented measurement.
Evaluation metrics should measure not just \emph{whether} an agent succeeds, but \emph{how efficiently} it succeeds and \emph{how} it fails.

\textbf{Outcome metrics.}
We define three primary outcome metrics.
\emph{Mean Reward (MR)} is the weighted average of task scores with resurrection bonus: $\text{MR} = \sum s_i w_i b_i / \sum w_i b_i$, where $w_i = [0.05, 0.20, 0.20, 0.15, 0.30, 0.10]$ and $b_i = 1.2$ (no resurrection) or $1.0$ (resurrected).
\emph{Pass Score (PS)} is the weighted pass rate: $\text{PS} = \sum p_i b_i^{\text{pass}} / (6 \times 1.5) \times 100$, where $p_i = 1$ if $s_i \geq 60$ and $b_i^{\text{pass}} = 1.5$ (no resurrection) or $1.0$.
\emph{Pipeline Score (PL)} is the simple average: $\text{PL} = \frac{1}{6}\sum s_i$.

\textbf{Process metrics.}
Beyond outcome scores, we record token consumption, interaction turns, command executions, retry counts, and elapsed wall-clock time per task and total.

\textbf{Failure taxonomy.}
We categorize agent failures into five types:
(1)~\emph{predecessor failure}---downstream unreachable due to missing artifact;
(2)~\emph{task-local reasoning failure}---cannot solve despite correct prerequisites;
(3)~\emph{execution failure}---build/dependency errors;
(4)~\emph{process collapse}---max-turn exhaustion, search loops, repetitive debugging without progress;
(5)~\emph{cost overrun}---budget exhausted.

This taxonomy enables systematic analysis of \emph{how} agents fail, not just \emph{whether} they fail.
For instance, process collapse is more prevalent in framework-based agents with unlimited iteration budgets, suggesting that scaffold design affects failure mode distribution.

\subsection{Models and Agent Backends}

We evaluate \eva{} across 22 model configurations spanning both proprietary and open-weight LLMs, using four agent backends:

\begin{table*}[t]
\centering
\small
\begin{tabular}{lll}
\toprule
\textbf{Backend} & \textbf{Type} & \textbf{Key Properties} \\
\midrule
OpenHands SDK & Full agent framework & Unlimited self-loop, skill injection, MCP \\
Codex CLI & CLI agent & Autonomous coding with sandboxed execution \\
Claude Code & CLI agent & File I/O, command execution, multi-turn \\
Kimi CLI & CLI agent & Conversational coding with plan management \\
\bottomrule
\end{tabular}
\caption{Agent backends used in \eva{} evaluation.}
\label{tab:backends}
\end{table}

All models are accessed through the \textsc{AIhub} unified API gateway, ensuring consistent infrastructure conditions across evaluations.

\subsection{Evaluation Settings}

We conduct evaluation under two primary settings:

\textbf{Setting~1: Serial pipeline with resurrection (default).}
Agents proceed through Tasks~0--5 in sequence.
When a task fails, resurrection is triggered and the agent continues with the golden intermediate artifact.
This setting produces the full leaderboard with both outcome and process metrics.

\textbf{Setting~2: Serial pipeline without resurrection.}
Agents proceed through Tasks~0--5 in sequence, but \emph{no resurrection is triggered}.
When a task fails, all subsequent tasks receive the agent's (possibly broken) artifact and typically cascade to zero scores.
This setting measures \emph{zero-shot pipeline depth}---how far an agent can progress without any external correction.

The comparison between Setting~1 and Setting~2 directly tests whether resurrection recovers diagnostic signal lost to cascading failure.

\subsection{Workload Comparison: \yatcc{} vs. \yatcc{}-Hard}

We evaluate models on both workloads to quantify the difficulty gap between scaffolded and from-scratch implementation.
\yatcc{} provides starter code and implementation frameworks; \yatcc{}-Hard removes all code logic, requiring agents to construct the full compiler from minimal specification.
The performance delta between the two workloads measures the value of scaffolding and reveals which models depend on provided structure versus which can operate autonomously.

Experiments are conducted under unified infrastructure conditions using identical execution environments and workload distributions, ensuring that observed differences reflect genuine model and scaffold capability gaps rather than environmental variability.